% TOP_PROBLEM: {"id": 30003917, "title": "Mass–Angular-Momentum Inequalities with Trapped and Multiple Horizons", "category_id": 16, "category": {"id": 16, "name": "physics", "display_name": "Mathematical Physics", "description": "Problems at the intersection of mathematics and physics.", "slug": "physics", "order_index": 16, "created_at": "2026-07-31T15:26:25.671Z"}, "status": "open"}
% FIRST_POSTED: 2026-10-06
% ATTEMPT_STATUS: unresolved
% OPENAI_CLAIM: {"schema_version":1,"source_repo":"https://github.com/openai/math","source_commit":"adc7f1241b42e322a6451854ab7e4b4c146bf78a","release_date":"2026-10-06","match":"partial","resolution":"partial","independently_reviewed":false,"second_pass":"n/a","families":[{"family":"260","title":"Spacetime Penrose inequalities: enclosing area, charge, rotation, and anti-de Sitter extensions","subject":"Mathematical physics","manuscripts":[{"title":"The Kerr–Newman Penrose Inequality for Axisymmetric Electrovacuum Exteriors","pdf_path":"preprints/The-Kerr-Newman-Penrose-Inequality-for-Axisymmetric-Electrovacuum-Exteriors-October-5-2026/paper.pdf","readme_path":"preprints/The-Kerr-Newman-Penrose-Inequality-for-Axisymmetric-Electrovacuum-Exteriors-October-5-2026/README.md","date":"2026-10-05","theorem_ref":"Theorem 2.2 (m:main) (main.tex:293)","statement_sources":[{"path":"preprints/The-Kerr-Newman-Penrose-Inequality-for-Axisymmetric-Electrovacuum-Exteriors-October-5-2026/build/main.tex","line":293,"label":"Theorem 2.2 (m:main)"}]}],"lean":{"doc_path":"lean/docs/260.md","comparators":[],"declarations":[],"covers_main_theorem":false},"reasoning_trace":null}]}
\documentclass[11pt]{article}
\usepackage[margin=25mm]{geometry}
\usepackage{fontspec}
\setmainfont{Latin Modern Roman}
\usepackage{amsmath,amssymb,mathtools}
\usepackage{unicode-math}
\setmathfont{Latin Modern Math}
\usepackage{enumitem}
\usepackage{xurl,hyperref,bookmark}
\hypersetup{colorlinks=true,urlcolor=blue}
\setcounter{secnumdepth}{0}
\setlength{\parindent}{0pt}
\setlength{\parskip}{5pt}
\setlength{\emergencystretch}{3em}
\begin{document}
\section*{Mass–Angular-Momentum Inequalities with Trapped and Multiple Horizons}

\subsection{OpenAI's claim}
OpenAI presents a mathematical result relevant to this question in the selected manuscript. For the designated axisymmetric electrovacuum class, the claimed bound is $m^2\geq A/(16\pi)+Q^2/2+\pi(Q^4+4J^2)/A$ with conserved total angular momentum.

\subsection{Mathematical statement of the claimed result}
\paragraph{Result 1.} Let the smooth oriented connected exterior be diffeomorphic to \(\mathbb R^3\) minus an open ball, with spherical smooth compact boundary \(S\), complete metric including \(S\), a single AF end, and smooth \(g,K,\mathcal E,\mathcal B\). A smooth effective \(U(1)\) action preserves all data and the boundary, with rotation generator of period \(2\pi\) agreeing with ordinary axial rotation at infinity. Impose source-free electrovacuum constraints
\[
16\pi\mu=R_g+(\operatorname{tr}_gK)^2-|K|^2,
\quad 8\pi J_{\rm con}=\operatorname{div}_g(K-(\operatorname{tr}_gK)g),
\]
\[
8\pi\mu=|\mathcal E|^2+|\mathcal B|^2,\quad
4\pi J_{\rm con}=(\mathcal E\times\mathcal B)^\flat,\quad
\operatorname{div}_g\mathcal E=\operatorname{div}_g\mathcal B=0.
\]
Use \(g-\delta=O_4(r^{-1})\), \(K=O_3(r^{-3})\), Coulomb leading fields
\(\mathcal E^i=Q_e x^i/r^3+O_3(r^{-3})\), \(\mathcal B^i=Q_b x^i/r^3+O_3(r^{-3})\), integrable constraints and finite ADM energy, momentum, angular momentum and charge fluxes. Require \(P_{\rm ADM}=0\), \(m=E_{\rm ADM}>0\), \(Q^2=Q_e^2+Q_b^2\), and let \(J\) denote conserved angular momentum, with its electromagnetic boundary contribution included; it agrees with \(J_{\rm ADM}\). The boundary is a future MOTS; no wholly exterior enclosing smooth surface is weakly outer trapped; \(\operatorname{Area}_g(\Gamma)\ge A=|S|_g>0\) for every full cut \(\Gamma\); and
\[A\ge4\pi\sqrt{Q^4+4J^2}.\]
\[
m^2\ge\frac{A}{16\pi}+\frac{Q^2}{2}+\frac{\pi(Q^4+4J^2)}{A}.
\]
Equality on the strict area branch \(A>4\pi\sqrt{Q^4+4J^2}\) identifies the entire original data with a smooth spacelike slice of subextremal dyonic Kerr--Newman with matching mass, angular momentum and charges, including smooth attachment and future normal at the horizon and an end at spatial infinity. Conversely, admissible such slices attain equality. The numerical endpoint is included, but extremal equality is not classified. Here every full cut is the entire boundary of an end-side exterior region, including all boundary-coincident components, and $O_j$ includes coordinate derivatives through order $j$.

\subsection{Outline of the argument}
Axial variables are compared with a Kerr--Newman family at prescribed conformal capacity. Riemannian scale comparison and electromagnetic angular-momentum bookkeeping give the mass bound; saturation on the strict physical-area branch reconstructs an admissible dyonic spacetime slice.

\subsection{Scope relative to this problem}
Optimizing the main Kerr–Newman area bound over A yields m\textasciicircum{}2>=|J| (and its charged refinement) in the manuscript axisymmetric connected electrovacuum horizon class. The problem statement also requests trapped inner boundaries generally and the multi-black-hole mass functional and minimizer, which are not settled. This entry records partial support; complete coverage of the question is not established by this catalogue.

\subsection{Formal verification}
For family 260, the only selected family entry is the CKS end-replacement construction and Schwarzschild equality examples. A connected-boundary inequality is conditional on the flat Penrose input; the general Bondi or spacetime inequalities are not formalized there. No Comparator is assigned to the selected manuscript here; the family page selects other papers or supporting statements. This site has not executed the Lean code. The scope descriptions do not certify the remaining manuscript arguments.

\subsection{Status caveat}
OpenAI attributes these manuscripts to work by an internal model. Its release cautions that arguments without formalization can contain mistakes. This catalogue has not independently checked the proofs; the labels report the mathematical claims attributed to OpenAI.

\subsection{Sources}
\begin{itemize}
\item OpenAI, \emph{The Kerr–Newman Penrose Inequality for Axisymmetric Electrovacuum Exteriors}, 2026-10-05; citation key \texttt{OAI:The-Kerr-Newman-Penrose-Inequality-for-Axisymmetric-Electrovacuum-Exteriors-October-5-2026}. \href{https://github.com/openai/math/blob/adc7f1241b42e322a6451854ab7e4b4c146bf78a/preprints/The-Kerr-Newman-Penrose-Inequality-for-Axisymmetric-Electrovacuum-Exteriors-October-5-2026/paper.pdf}{PDF (GitHub)}; \href{https://github.com/openai/math/raw/adc7f1241b42e322a6451854ab7e4b4c146bf78a/preprints/The-Kerr-Newman-Penrose-Inequality-for-Axisymmetric-Electrovacuum-Exteriors-October-5-2026/paper.pdf}{PDF (direct)}; \href{https://github.com/openai/math/blob/adc7f1241b42e322a6451854ab7e4b4c146bf78a/preprints/The-Kerr-Newman-Penrose-Inequality-for-Axisymmetric-Electrovacuum-Exteriors-October-5-2026/README.md}{Manuscript page and citation}; \href{https://github.com/openai/math/blob/main/preprints/The-Kerr-Newman-Penrose-Inequality-for-Axisymmetric-Electrovacuum-Exteriors-October-5-2026/paper.pdf}{Latest version}; Statement source: \href{https://github.com/openai/math/blob/adc7f1241b42e322a6451854ab7e4b4c146bf78a/preprints/The-Kerr-Newman-Penrose-Inequality-for-Axisymmetric-Electrovacuum-Exteriors-October-5-2026/build/main.tex\#L293}{Theorem 2.2 (m:main)}.
\item Family 260: \href{https://github.com/openai/math/blob/adc7f1241b42e322a6451854ab7e4b4c146bf78a/lean/docs/260.md}{Lean scope}.
\item \href{https://github.com/openai/math/blob/adc7f1241b42e322a6451854ab7e4b4c146bf78a/CONTENTS.md}{OpenAI catalogue entry}.
\end{itemize}
\end{document}
